\documentclass[12pt,a4paper]{ctexart}

% 页面设置
\usepackage[
    left=2.5cm,
    right=2.5cm,
    top=2.3cm,
    bottom=2.3cm
]{geometry}

% 常用宏包
\usepackage{xcolor}
\usepackage{tcolorbox}
\usepackage{titlesec}
\usepackage{fancyhdr}
\usepackage{enumitem}
\usepackage{tabularx}
\usepackage{amsmath, amssymb}
\usepackage{hyperref}

\tcbuselibrary{skins, breakable}

% ======================
% 颜色主题：清爽学习笔记风
% ======================

% 主色：深蓝绿
\definecolor{MainColor}{HTML}{2A6F73}

% 辅助色：青绿色
\definecolor{SecondColor}{HTML}{3CAEA3}

% 浅背景色
\definecolor{LightMain}{HTML}{EAF7F6}

% 标题深色
\definecolor{TitleDark}{HTML}{1F3A40}

% 文字灰
\definecolor{TextGray}{HTML}{555555}

% 边框灰
\definecolor{BorderGray}{HTML}{CBD5D6}

% 提醒色：柔和橙
\definecolor{WarnColor}{HTML}{F2B35D}

% 易错点红色
\definecolor{ErrorColor}{HTML}{C94C4C}

% 定义框紫色
\definecolor{DefineColor}{HTML}{6C63B7}

% ======================
% 正文设置
% ======================

\linespread{1.25}
\setlength{\parindent}{2em}
\setlength{\parskip}{0.25em}

\hypersetup{
  colorlinks=true,
  linkcolor=MainColor,
  urlcolor=MainColor
}

% ======================
% 页眉页脚
% ======================

\pagestyle{fancy}
\fancyhf{}
\fancyfoot[L]{\textcolor{MainColor}{机器学习笔记}}
\fancyfoot[R]{\textcolor{MainColor}{\thepage}}
\renewcommand{\headrulewidth}{0pt}
\renewcommand{\footrulewidth}{0pt}

% ======================
% 顶部小标题栏
% ======================

\newcommand{\topbar}[2]{
  \noindent
  {\small\textcolor{MainColor}{#1}}
  \hfill
  {\small\bfseries\textcolor{TitleDark}{#2}}
  \par\vspace{0.35em}
  {\color{BorderGray}\hrule}
  \vspace{1.2em}
}

% ======================
% 标题卡片
% ======================

\newcommand{\notetitlecard}[6]{
\begin{tcolorbox}[
  enhanced,
  colback=white,
  colframe=BorderGray,
  boxrule=0.6pt,
  arc=2mm,
  left=12pt,
  right=12pt,
  top=10pt,
  bottom=10pt,
  borderline west={4pt}{0pt}{SecondColor}
]
\begin{tabularx}{\linewidth}{X r}
{\small\textcolor{TextGray}{#1}} & {\small\textcolor{TextGray}{#2}} \\[0.9em]
\multicolumn{2}{l}{
  {\Huge\bfseries\textcolor{MainColor}{#3}}
} \\[0.45em]
\multicolumn{2}{l}{
  {\large\textcolor{SecondColor!90!black}{#4}}
} \\[0.45em]
\multicolumn{2}{l}{
  {\small\textcolor{TextGray}{课程：#5 \qquad 作者：#6}}
}
\end{tabularx}
\end{tcolorbox}
\vspace{1em}
}

% ======================
% 摘要框
% ======================

\newtcolorbox{summarybox}[1][本节摘要]{
  enhanced,
  breakable,
  colback=LightMain,
  colframe=SecondColor,
  colbacktitle=SecondColor,
  coltitle=white,
  title={#1},
  fonttitle=\bfseries,
  boxrule=0.6pt,
  arc=2mm,
  left=12pt,
  right=12pt,
  top=10pt,
  bottom=10pt
}

% ======================
% 定义框
% ======================

\newtcolorbox{definitionbox}[1][定义]{
  enhanced,
  breakable,
  colback=DefineColor!6,
  colframe=DefineColor,
  colbacktitle=DefineColor,
  coltitle=white,
  title={#1},
  fonttitle=\bfseries,
  boxrule=0.6pt,
  arc=2mm,
  left=12pt,
  right=12pt,
  top=10pt,
  bottom=10pt
}

% ======================
% 重点框
% ======================

\newtcolorbox{keybox}[1][重点]{
  enhanced,
  breakable,
  colback=MainColor!5,
  colframe=MainColor,
  colbacktitle=MainColor,
  coltitle=white,
  title={#1},
  fonttitle=\bfseries,
  boxrule=0.6pt,
  arc=2mm,
  left=12pt,
  right=12pt,
  top=10pt,
  bottom=10pt
}

% ======================
% 例题框
% ======================

\newtcolorbox{examplebox}[1][例题]{
  enhanced,
  breakable,
  colback=white,
  colframe=SecondColor,
  colbacktitle=SecondColor!85!black,
  coltitle=white,
  title={#1},
  fonttitle=\bfseries,
  boxrule=0.6pt,
  arc=2mm,
  left=12pt,
  right=12pt,
  top=10pt,
  bottom=10pt
}

% ======================
% 易错点框
% ======================

\newtcolorbox{mistakebox}[1][易错点]{
  enhanced,
  breakable,
  colback=ErrorColor!5,
  colframe=ErrorColor,
  colbacktitle=ErrorColor,
  coltitle=white,
  title={#1},
  fonttitle=\bfseries,
  boxrule=0.6pt,
  arc=2mm,
  left=12pt,
  right=12pt,
  top=10pt,
  bottom=10pt
}

% ======================
% 提醒框
% ======================

\newtcolorbox{tipbox}[1][提示]{
  enhanced,
  breakable,
  colback=WarnColor!10,
  colframe=WarnColor,
  colbacktitle=WarnColor,
  coltitle=white,
  title={#1},
  fonttitle=\bfseries,
  boxrule=0.6pt,
  arc=2mm,
  left=12pt,
  right=12pt,
  top=10pt,
  bottom=10pt
}

% ======================
% 章节标题样式
% ======================

\titleformat{\section}
  {\Large\bfseries\color{MainColor}}
  {\thesection}
  {0.8em}
  {}

\titleformat{\subsection}
  {\large\bfseries\color{TitleDark}}
  {\thesubsection}
  {0.8em}
  {}

\titleformat{\subsubsection}
  {\normalsize\bfseries\color{SecondColor!80!black}}
  {\thesubsubsection}
  {0.8em}
  {}

% ======================
% 列表样式
% ======================

\setlist[itemize]{
  leftmargin=2.2em,
  itemsep=0.3em,
  topsep=0.4em
}

\setlist[enumerate]{
  leftmargin=2.2em,
  itemsep=0.3em,
  topsep=0.4em
}

% ======================
% 自定义强调命令
% ======================

\newcommand{\important}[1]{\textbf{\textcolor{MainColor}{#1}}}
\newcommand{\warning}[1]{\textbf{\textcolor{ErrorColor}{#1}}}
\newcommand{\concept}[1]{\textbf{\textcolor{DefineColor}{#1}}}

% ======================
% 正文开始
% ======================

\begin{document}

\topbar{吴恩达机器学习}{学习笔记}

\notetitlecard
  {Study Notes Series}
  {2026 Spring}
  {吴恩达机器学习笔记(2)}
  {}
  {机器学习}
  {C++和草稿纸}

\begin{summarybox}
本节介绍多变量线性回归的向量化表示、梯度下降与特征缩放，并总结收敛诊断、特征工程和多项式回归。
\end{summarybox}
\section{多变量线性回归}

\subsection{多维特征}

\begin{definitionbox}[多维特征的符号表示]
当一个训练样本包含多个输入变量时，这些输入变量称为\concept{特征（features）}。设每个样本共有 $n$ 个特征，则第 $i$ 个训练样本可写成一个 $n$ 维列向量：
\[
  \mathbf{x}^{(i)}=
  \begin{bmatrix}
    x_1^{(i)} & x_2^{(i)} & \cdots & x_n^{(i)}
  \end{bmatrix}^{T}.
\]
其中，$x_j^{(i)}$ 表示第 $i$ 个训练样本的第 $j$ 个特征。例如，若
$\mathbf{x}^{(2)}=[1416,3,2,40]^T$，则 $x_2^{(2)}=3$，$x_3^{(2)}=2$。

具有 $n$ 个特征的线性回归假设函数为
\[
  h_{\boldsymbol{\theta}}(\mathbf{x})
  =\theta_0+\theta_1x_1+\theta_2x_2+\cdots+\theta_nx_n.
\]
令 $x_0=1$，并记
$\mathbf{x}=[x_0,x_1,\ldots,x_n]^T$、
$\boldsymbol{\theta}=[\theta_0,\theta_1,\ldots,\theta_n]^T$，则可将假设函数简写为
\[
  h_{\boldsymbol{\theta}}(\mathbf{x})=\boldsymbol{\theta}^{T}\mathbf{x}.
\]
此时每个训练样本和参数向量都是 $(n+1)$ 维向量；若共有 $m$ 个训练样本，则特征矩阵 $X$ 的维度为 $m\times(n+1)$。
\end{definitionbox}

\begin{keybox}[向量化的好处]
对于含有多个特征的模型，若逐项计算预测值，需要依次完成求和：
\[
  h_{\boldsymbol{\theta}}(\mathbf{x})
  =\sum_{j=0}^{n}\theta_jx_j.
\]
使用向量化后，同一计算可直接写成点积：
\[
  h_{\boldsymbol{\theta}}(\mathbf{x})
  =\boldsymbol{\theta}^{T}\mathbf{x}.
\]
向量化主要有以下好处：
\begin{itemize}
  \item \textbf{代码更简洁：}无需显式编写循环，模型表达式更接近数学公式，也更不容易出现索引错误。
  \item \textbf{计算速度更快：}NumPy 等数值计算库会使用经过优化的底层实现，比在高级语言中逐项运算更高效。
  \item \textbf{能够并行计算：}点积和矩阵运算可以利用 CPU 的向量指令或 GPU 的并行能力，同时处理多个元素。
\end{itemize}
当 $m$ 个训练样本组成特征矩阵 $X$ 时，还可以一次得到全部预测结果：
\[
  \hat{\mathbf{y}}=X\boldsymbol{\theta}.
\]
因此，随着特征数或训练样本数增加，向量化通常能带来更明显的效率优势。
\end{keybox}

\subsection{多变量梯度下降}

\begin{definitionbox}[多变量线性回归的代价函数]
多变量线性回归与单变量线性回归使用相同的平方误差代价函数。对于 $m$ 个训练样本，模型预测和代价函数分别为
\[
  h_{\boldsymbol{\theta}}\!\left(\mathbf{x}^{(i)}\right)
  =\boldsymbol{\theta}^{T}\mathbf{x}^{(i)},
\]
\[
  J(\boldsymbol{\theta})
  =\frac{1}{2m}\sum_{i=1}^{m}
  \left(h_{\boldsymbol{\theta}}\!\left(\mathbf{x}^{(i)}\right)-y^{(i)}\right)^2.
\]
训练的目标是找到一组参数 $\boldsymbol{\theta}$，使 $J(\boldsymbol{\theta})$ 尽可能小。
\end{definitionbox}

\begin{keybox}[多变量的批量梯度下降]
对每个参数 $\theta_j$，梯度下降的更新规则为
\[
  \theta_j := \theta_j
  -\alpha\frac{\partial}{\partial\theta_j}J(\boldsymbol{\theta})
  =\theta_j-\alpha\frac{1}{m}\sum_{i=1}^{m}
  \left(h_{\boldsymbol{\theta}}\!\left(\mathbf{x}^{(i)}\right)-y^{(i)}\right)x_j^{(i)},
  \qquad j=0,1,\ldots,n.
\]
其中，$\alpha$ 是学习率。由于 $x_0^{(i)}=1$，同一个公式也适用于偏置参数 $\theta_0$。

一次完整迭代可概括为：
\begin{enumerate}
  \item 使用当前参数计算所有训练样本的预测值与误差；
  \item 根据全部 $m$ 个训练样本计算每个参数的梯度；
  \item 使用更新前的同一组参数，\important{同时更新} $\theta_0,\theta_1,\ldots,\theta_n$；
  \item 重复以上过程，直到代价函数收敛。
\end{enumerate}
每次更新都使用整个训练集，因此这种方法称为\concept{批量梯度下降（Batch Gradient Descent）}。
\end{keybox}

\begin{tipbox}[向量化更新]
令 $\mathbf{y}=[y^{(1)},y^{(2)},\ldots,y^{(m)}]^T$，则全部样本的梯度和参数更新可以写成
\[
  \nabla J(\boldsymbol{\theta})
  =\frac{1}{m}X^T\left(X\boldsymbol{\theta}-\mathbf{y}\right),
  \qquad
  \boldsymbol{\theta}:=\boldsymbol{\theta}
  -\alpha\nabla J(\boldsymbol{\theta}).
\]
该形式一次计算所有参数的梯度，代码更简洁，也能利用数值计算库的矩阵运算加速。
\end{tipbox}
\subsection{梯度下降法的一些实践}

\subsubsection{特征缩放}

\begin{keybox}[为什么需要特征缩放]
当不同特征的取值范围相差很大时，代价函数的等高线会变得狭长。梯度下降可能在较窄的方向上来回震荡，需要许多次迭代才能到达最小值。

\concept{特征缩放（Feature Scaling）}将各个特征调整到相近的数值范围，使代价函数的等高线更接近圆形。这样梯度下降可以更直接地朝最小值移动，通常能够使用更合适的学习率并更快收敛。缩放后的特征不必严格落在同一区间，重点是避免不同特征之间存在悬殊的数量级差异。
\end{keybox}

\begin{definitionbox}[均值归一化（Mean Normalization）]
对第 $j$ 个特征，记训练集中的均值为 $\mu_j$，最大值和最小值分别为 $x_{j,\max}$ 与 $x_{j,\min}$。均值归一化定义为
\[
  x_j \leftarrow
  \frac{x_j-\mu_j}{x_{j,\max}-x_{j,\min}}.
\]
减去均值会使特征以 $0$ 为中心，再除以取值范围则缩小其尺度。处理后的数值通常位于 $[-1,1]$ 附近。
\end{definitionbox}

\begin{definitionbox}[Z-score 标准化（Z-score Normalization）]
先使用训练集计算第 $j$ 个特征的均值和标准差：
\[
  \mu_j=\frac{1}{m}\sum_{i=1}^{m}x_j^{(i)},
  \qquad
  \sigma_j=\sqrt{\frac{1}{m}\sum_{i=1}^{m}
  \left(x_j^{(i)}-\mu_j\right)^2}.
\]
然后进行 Z-score 标准化：
\[
  x_j \leftarrow \frac{x_j-\mu_j}{\sigma_j}.
\]
标准化后的特征均值为 $0$、标准差为 $1$。它不保证所有数据都位于 $[-1,1]$ 内，但通常能让不同特征具有相近的尺度。
\end{definitionbox}

\begin{tipbox}[使用特征缩放时的注意事项]
\begin{itemize}
  \item 均值、最大值、最小值和标准差只能根据\important{训练集}计算；验证集、测试集及新输入必须复用同一组统计量，以避免数据泄漏。
  \item 偏置特征 $x_0=1$ 不需要缩放。若某个特征的取值范围或标准差为 $0$，说明它在训练集中是常数，应移除该特征或单独处理。
  \item Mean normalization 和 Z-score normalization 都能用于特征缩放。选择其中一种并在所有数据上保持一致即可。
\end{itemize}
\end{tipbox}
\subsubsection{检查梯度下降是否收敛}

\begin{keybox}[通过学习曲线判断收敛]
在每次迭代后记录当前代价
$J^{(t)}=J\!\left(\boldsymbol{\theta}^{(t)}\right)$，并绘制横轴为迭代次数 $t$、纵轴为 $J^{(t)}$ 的\concept{学习曲线（Learning Curve）}。

若梯度下降实现正确且学习率合适，$J^{(t)}$ 应当在每次迭代后下降。曲线逐渐变平且相邻代价变化很小时，算法可能已接近最小值。所需迭代次数因问题而异，不能只靠预设次数判断收敛。
\end{keybox}

\begin{definitionbox}[自动收敛测试]
设第 $t$ 次迭代带来的代价下降量为
\[
  \Delta J^{(t)}=J^{(t-1)}-J^{(t)}.
\]
选取一个很小的阈值 $\varepsilon>0$。当
\[
  0\leq \Delta J^{(t)}\leq\varepsilon
\]
时，可认为本次迭代对代价函数的改善已经很小，并据此判断梯度下降已经收敛。吴恩达老师以 $\varepsilon=10^{-3}$ 为例，但该数值不是固定标准，需要根据代价函数的尺度和具体任务选择。

实际应用中，$\varepsilon$ 往往不容易设定，因此先观察学习曲线通常是更直观的检查方式；需要自动停止训练时，再配合收敛阈值使用。
\end{definitionbox}
\subsubsection{学习率的选择}
\begin{mistakebox}[学习曲线异常时如何排查]
\begin{itemize}
  \item 若 $J^{(t)}$ 随迭代\warning{持续上升}，通常表示学习率 $\alpha$ 过大，或者代价函数、梯度及参数更新的实现存在错误。
  \item 若曲线明显上下震荡，梯度下降可能在最小值两侧反复跨越。应尝试减小 $\alpha$，并检查各个特征是否已经缩放到相近范围。
  \item 若 $J^{(t)}$ 始终下降但速度很慢，可以适当增大 $\alpha$，同时通过特征缩放改善收敛速度。
  \item 只有当代价函数正常下降并逐渐变平时，才能将“变化量很小”解释为收敛；代价上升不能被自动停止条件误判为收敛。
\end{itemize}
\end{mistakebox}
\subsubsection{特征工程}

\begin{definitionbox}[什么是特征工程]
\concept{特征工程（Feature Engineering）}是利用对问题的理解，将已有特征进行选择、组合或变换，从而构造出更能表达数据规律的新特征。合适的新特征可以让较简单的模型表示原本难以拟合的关系。

特征工程改变的是输入 $\mathbf{x}$ 的表示方式，而线性回归模型对参数 $\boldsymbol{\theta}$ 仍然保持线性，因此仍可使用相同的代价函数和梯度下降算法训练。
\end{definitionbox}

\begin{examplebox}[房价预测中的组合特征]
假设房屋的临街宽度为 $x_1=\mathrm{frontage}$，纵向深度为 $x_2=\mathrm{depth}$。只使用原始特征时，模型为
\[
  h_{\boldsymbol{\theta}}(\mathbf{x})
  =\theta_0+\theta_1x_1+\theta_2x_2.
\]
根据房价与占地面积相关的领域知识，可以构造新特征
\[
  x_3=x_1x_2
  =\mathrm{frontage}\times\mathrm{depth}
  =\mathrm{area},
\]
并将模型扩展为
\[
  h_{\boldsymbol{\theta}}(\mathbf{x})
  =\theta_0+\theta_1x_1+\theta_2x_2+\theta_3x_3.
\]
乘积特征 $x_1x_2$ 表达了两个原始特征之间的交互关系，也使模型能够直接利用“面积”这一更有意义的概念。
\end{examplebox}

\begin{tipbox}[构造特征的原则]
新特征应当有明确的实际意义或能够表达观察到的数据形状，例如乘积、比值、平方项或平方根等。特征并非越多越好，应通过训练效果和验证集表现判断新特征是否真正有用。
\end{tipbox}

\subsubsection{多项式回归}

\begin{definitionbox}[多项式回归]
当数据呈现明显的曲线趋势时，一条直线可能无法很好地拟合数据。可以由原始特征 $x$ 构造
\[
  x_1=x,\qquad x_2=x^2,\qquad x_3=x^3,
\]
再使用线性回归得到三次多项式模型
\[
  h_{\boldsymbol{\theta}}(x)
  =\theta_0+\theta_1x+\theta_2x^2+\theta_3x^3.
\]
模型虽然对输入 $x$ 是非线性的，但对参数 $\theta_0,\theta_1,\theta_2,\theta_3$ 仍然是线性的，所以仍属于线性回归，并可沿用原来的训练方法。
\end{definitionbox}

\begin{keybox}[选择合适的特征变换]
不同的特征变换可以产生不同形状的拟合曲线，例如
\[
  \text{二次：}\quad
  h_{\boldsymbol{\theta}}(x)=\theta_0+\theta_1x+\theta_2x^2,
  \qquad
  \text{平方根：}\quad
  h_{\boldsymbol{\theta}}(x)=\theta_0+\theta_1x+\theta_2\sqrt{x}.
\]
应先观察数据的整体趋势，再选择二次项、三次项、平方根或其他有意义的变换，并使用验证集比较不同模型，而不是只凭训练误差决定模型形式。
\end{keybox}

\begin{mistakebox}[多项式回归的注意事项]
\begin{itemize}
  \item 幂次会迅速放大特征尺度。例如 $x$、$x^2$ 和 $x^3$ 的取值范围可能相差很多，因此使用梯度下降前\important{必须对各个多项式特征进行缩放}。
  \item 多项式次数过低可能无法表达数据规律；次数过高则可能跟随训练数据中的噪声，导致过拟合。
  \item 高次多项式在训练数据范围之外可能产生不合理的快速增长或下降，因此需要谨慎解释外推结果。
\end{itemize}
\end{mistakebox}

\end{document}
